\section{Sustainability}
\label{sec:sustainability}

\subsection{Self-Assessment of the Sustainability Survey}
\label{subsec:self-assessment-sustainability-survey}

I was aware that modern AI is highly resource-intensive, but quantifying the exact electricity draw and operational cost of my local GPUs makes the issue undeniable. It is one thing to read articles about the massive environmental impact of cloud data centers, but it hits differently when you are measuring the watts your own pipeline is burning. 

Rather than feeling discouraged, this realization gives me a strong sense of purpose. The efficiency frameworks I am building are not just theoretical math exercises; they are direct tools for sustainability. By successfully compressing the information bottleneck and running powerful reasoning models locally on consumer-grade hardware, we can drastically reduce the reliance on energy-hungry cloud APIs. 

However, the survey exposed a gap in my perspective. While I am confident managing the economic metrics of the project—tracking CAPEX returns, hardware amortization, and labor rates—I lack the frameworks to fully measure broader social impacts, such as how deploying these enterprise AI tools might disrupt local IT jobs. Recognizing this blind spot ensures this project remains holistic, reminding me that elite technical performance must always be paired with socio-environmental responsibility.

\subsection{Economic Impact}
\label{subsec:economic-impact}

\subsubsection{Economic dimension of the sustainability matrix}
\label{subsubsec:economic-dimension-sustainability-matrix}

By applying corporate rigor, real-world accounting, fully burdened labor rates, and Expected Monetary Value (EMV) risk provisioning, this project functions as a genuine enterprise contract. This meticulous planning ensures operational viability and tracks every euro to justify the investment.

\subsubsection{How does the solution improve economic issues with respect to other existing solutions?}
\label{subsubsec:solution-improves-economic-issues}

Industry standards rely on renting cloud GPUs or per-token APIs, trapping companies in unpredictable operational expenses (OPEX). By mathematically optimizing the context window for high-fidelity local reasoning, this solution shifts to a one-time CAPEX hardware investment. This achieves a four-month break-even point, drastically lowering the financial barrier for smaller businesses.

\subsection{Environmental Impact}
\label{subsec:environmental-impact}

\subsubsection{Have you estimated the environmental impact of the project?}
\label{subsubsec:estimated-environmental-impact}

Moving beyond abstract concepts, hardware telemetry shows 540 nominal hours ($0.30 \text{ kW}$) and 300 training hours ($0.85 \text{ kW}$) consume $417 \text{ kWh}$. At the Spanish grid's emission factor ($\approx 0.20 \text{ kg CO}_2/\text{kWh}$), this yields $83.4 \text{ kg CO}_2$. Furthermore, 90 urban commutes (882 km) in a 7.2 L/100 km vehicle consume 63.5 liters of fuel, generating $146.7 \text{ kg CO}_2$. The total footprint is $230.1 \text{ kg CO}_2$, proving that even software exacts a physical toll.

\subsubsection{Did you plan to minimize its impact, for example, by reusing resources?}
\label{subsubsec:minimize-environmental-impact}

To minimize this footprint, I am utilizing an existing industrial workspace and a hardware stack that will be fully repurposed for long-term corporate R\&D at \textit{Decolletatge Farrés S.A.} post-defense. At the software level, optimizing the C++ solver to prune irrelevant data drastically reduces total execution time. This algorithmic efficiency ensures the GPUs spend significantly less time drawing maximum wattage during inference, directly lowering the operational carbon footprint.

\subsubsection{How will your solution improve the environment with respect to other existing solutions?}
\label{subsubsec:improve-environmental-impact}

Unlike standard AI paradigms that brute-force massive models in energy-heavy, centralized data centers, my solution uses game-theoretic synergy to compress the information bottleneck. This allows complex, multi-hop reasoning to run efficiently on standard consumer-grade hardware. By eliminating the need for continuous cloud data transmission and massive parallel computing clusters, this localized approach achieves elite AI performance without contributing to the escalating carbon and water footprint of industrial server farms.

\subsection{Social Impact}
\label{subsec:social-impact}

\subsubsection{What do you think you will achieve – in terms of personal growth – from doing this project?}
\label{subsubsec:personal-growth}

Acting as the bridge between my academic studies and my future professional career, building this architecture tested my resilience against deep algorithmic bottlenecks, hardware constraints, and strict project deadlines. More importantly, taking on the role of project manager forced me to understand how decisions intersect with financial and social realities.

\subsubsection{How will your solution improve the quality of life with respect to other existing solutions?}
\label{subsubsec:improve-quality-life}

Currently, the most advanced AI tools are locked behind expensive, proprietary APIs. By making highly efficient reasoning viable on locally hosted, accessible hardware, this solution actively democratizes AI. It levels the technological playing field, empowering local businesses, independent researchers, and smaller engineering teams to automate complex analytical tasks and remain highly competitive without being priced out by multinational tech monopolies.

\subsubsection{Is there a real need for the project?}
\label{subsubsec:real-need-project}

There is an urgent need for this architecture. Sectors such as legal compliance, healthcare, and proprietary corporate R\&D handle highly sensitive data that, under strict frameworks like the GDPR, legally cannot be transmitted to third-party cloud providers. These industries require intelligent, entirely offline reasoning systems. By solving the memory and latency bottlenecks that currently prevent local execution, this project provides the secure, private, and financially accessible offline alternative that the enterprise market severely lacks today.